\chapter{Related Work and theory}
This thesis is inspired from the delivery methods utilized by Content delivery
Networks (CDNs) \citep{cdn3} \citep{cdn4} \citep{cdn5} \citep{cdn6} and torrent networks in order to relay large files on a network in an
efficient manner. CDNs transmit large file over an adaptive relay network in the
form of software or media distribution. In order to transmit these large files
to multiple thousands of clients, the CDN caches symbols in its memory and opens
up many fountains one for each client.\citep{cdn1} and \citep{cdn2} show
efficient ways of using fountain codes to deliver a continuous stream
of systematic symbols (Systematic symbols : The first k that are transmitted are the encoded
data part). Assuming that a part of those k symbols are dropped, repair symbols
are sent till the client collects a total of k ( original + repair) symbols and
is able to reconstruct this data.

Torrents on the other hand use a distributed file in order to find the nodes
that can provide the information. This file is called the torrent file. When we
know the nodes that can provide the information, we request them for the
information. These nodes respond with a part of the file that they own. When all
nodes provide the known information, the file is aggregated and reconstructed
into the usable file.

BitTorrents \citep{bittorrent1} \citep{cohem} were experimented to use fountain codes in 2012 by \citep{torrent-fountain}.
However, the results in the case of torrents were almost comparable to the one
with no coding technique used in the initial phase, but as enough seeders seed
the data, the download with coding increases to a much faster rate. In our case,
the block will be present with all the seeders, therefore the algorithm to
find/disperse data is not employed as is the case with this paper.

In the context of blockchains, miners who need to send the information to many
nodes and in a fast and reliable manner, stream multiple fountains for each node
that they are connected to. The full nodes use the torrent approach of
aggregating data and upload the partial information.

Blockchains, in the near future, in order to scale on chain will require to
transmit large blocks over the network in a relatively small block interval.
Since, CDNs effectively transmits their target data for efficient downloads, the
idea can be abstracted and specific pieces of technology can be utilized in
blockchains to synchronize large blocks over the decentralized distributed
network.

\section{Digital Fountains Codes}
The common analogy to understand Fountain codes, is a water fountain. In order to fill a
glass of water, one needs enough water, and not specific parts of water.
Fountain codes work on the same principle. In order to reconstruct data, one
needs enough symbols, here $k$. They can be a combination of source and repair
symbols.

Fountain codes are also known as rateless erasure codes. Near optimal erasure
codes \citep{efficient-erasure} can transform $k$ symbol messages into practically infinitely encoded form
i.e. they can generate an arbitrary amount of redundancy symbols that can be
used for error correction.
\textbf{Luby Transform} \citep{lt} codes are the first class of practical fountain codes
which depend on bipartite graphs to trade reception overhead for encoding
and decoding speed. 

These codes employ a simple $\oplus $ based algorithm to encode and decode a
message. And these codes adopt only a one way communication protocol which can
directly be plugged as an alternative to \texttt{getdata} request explained in
the coming chapters.

\subsection{Encoding and Decoding of Luby Transform Codes}
\begin{enumerate}
    \item The sender encodes and sends packets of information.
    \item The receiver evaluates each packet. If the packet has an error, it
        discards that packet, otherwise it is saved as a message.
    \item Eventually the receiver has enough valid packets to reconstruct the
        message, on which an ACK is sent after transmission.
\end{enumerate}

\subsubsection{Luby Transform Encoding Process}
\begin{enumerate}
    \item
        The message is divided into $n$ blocks of roughly equal length.
    \item
        The degree $ d, 1 <= d <= n $ is chosen at random.
    \item
        $ d$ blocks are chosen at random.
    \item
        If $M_i$ is the $i^th$ block of the message, data portion of the next
        packet is $M_i1 \oplus M_i2 \oplus ... \oplus M_id $
    \item
        A prefix is appended to the encoded packet: defining $n,d$ and a list of
        indices $i_1, i_2, .. i_d$.
    \item
        Some form of error detecting code (as simple as cyclic redundancy check)
        is applied to the packet and transmitted.
\end{enumerate}

\subsubsection{Luby Transform Decoding Process}
\begin{enumerate}
    \item The packet is first checked, if the packet isn't clean or it has
        already been sent, it is discarded.
    \item 
        If the cleanly received packet is of degree $d > 1$, it is first
        processed against all the fully decoded blocks in the message queuing
        area, then stored in a buffer area if its reduced degree is greater than
        1.
    \item
        If $d=1$ is received, the packet is matched against all the packets of
        $d>1$ residing in the buffer. If $d=2$ is received, it is reduced to
        $d=1$ and moved to message queuing area.
    \item When all the $n$ blocks are produced, the signal is transmitted.
\end{enumerate}


\subsection{Raptor Codes}
Raptor codes are systematic fountain codes, where on receiving symbols the
system can't identify whether these are original source or the repair symbols.

The way these are generated is, from the original source $k$ symbols are
constructed that can decode the entire source. Once we know that, the $k$ can
decode the original, we generate more symbols from the original symbols again
and call them the repair symbols. The receiver end now gets a degree
distribution in such a way that we don't actually require symbols of degree 1 to
decode all the original symbols.The degree distribution employed for RaptorQ is
taken from a mix of 2 or more distributions. Any mix of source and repair symbols, as
long as we have enough symbols ($k$), we can decode the source. 

The idea that makes Raptor codes work they do, is deactivation decoding. In
order to decode the entire data, we find a very small fraction of symbols
through Giant Component analysis and which can be later discovered using
Gaussian elimination as shown in \citep{raptor-dev} and \citep{raptorQ}, that if we assume to be decoded, we can decode the
rest of the data.

With this version improvised over time, if $k$ symbols are received at the
receiver's end, there is a $.5\%$ chance of failure and is independent of $k$, and with every  overhead
symbol, the chances of failure to decode reduces by a factor of 200 as stated in \citep{fountain-pres}. 
 
Therefore, for the current time, the research on erasure channels have reached
to optimality with the invention of raptorQ codes with a clever use of two or
more degree distributions.

The result on Raptor codes shows that a file on TCP that takes about 2min 15
seconds to downloads, downloads in about 13 seconds.
Therefore, this research is further applied to recover/reconstruct data in
blockchain in order to understand its implications on block propagation.

Some of the advantages of these codes, as stated in the paper \cite{multicast}

\begin{enumerate}
    \item \textbf{Scalable:} The server (miner) load remains constant if the
        number of leeching nodes increase. In cases of miners relaying blocks on
        asynchronous multicast connections, this will help increase the
        performance in sending data.
    \item \textbf{Reliable:} An exact copy of the original file is reconstructed
        at the receiver end, which is very important, since we are dealing with
        transactional data. 
    \item \textbf{Reception Efficient:} The total number of packets that each
        receiver needs to reconstruct  the file or data is minimal. In ideal
        cases, the length of data to be received equals the length of aggregated
        symbols. As we add more symbols, every symbols decreases the chance of
        decode failure by 200.
    \item \textbf{Time Efficient:}The time spent to generate and reconstruct the
        packets is minimal. These symbols can be saved in a file and can be
        directly broadcasted in order to respond to a request.
    \item \textbf{Time Independent:} Receivers may initiate the download at their
        discretion, implying that different receivers may start the download at
        widely disparate times. Receivers may sporadically be interrupted and
        continue the download at a later time.
    \item \textbf{Server Independent:} Receivers may collect packets for the
        file  from  one  or  more  servers  that  are  transmitting packets. No
        coordination between servers should be required for this.
    \item \textbf{Tolerant:} The protocol should tolerate a heterogeneous
        population of receivers, especially a variety of end-to-end packet loss
        rates and data rates.
\end{enumerate}

